% ============================================================
% STATUS (2026-06-25): TO BE CONTINUED / UNDER REVIEW.
% This section may be cut. Open question is whether the Dutch
% reference case earns its place or whether its content folds
% into the introduction/solution. Do not treat as final.
% Outstanding: measured Dutch greenhouse yield (kg/m2) still
% needed if section is kept (see sources/citation_map.md +
% memory project_netherlands_case_study).
% ============================================================
\section{Background: controlled-environment agriculture and the Dutch reference case}

Section~\ref{sec:problems} described a production model whose harms follow in large part from its exposure: crops grown in open fields across vast areas, subject to weather, pests and the erosion of the soil beneath them, and sustained by the chemical and water inputs that such exposure demands. One way to address harms that stem from exposure is to remove it. Controlled-environment agriculture (CEA) encloses cultivation within a structure, typically a greenhouse, in which temperature, humidity, irrigation and nutrient supply can be regulated and production is largely decoupled from local soil and weather. Several of the problems documented earlier follow from the open-field model directly, and enclosure addresses them at their source: it can eliminate tillage-driven soil erosion, recirculate water and nutrients instead of leaching them into waterways, and exclude pests without broadcast spraying. This is why CEA is the point of departure for the system proposed in this paper, the open-source and decentralised design developed in the sections that follow.

It is worth establishing first that the approach is not speculative. High-tech greenhouse production has been practised at national scale for decades, and the Netherlands is its clearest example. The Dutch sector is useful here in two ways. It demonstrates empirically what the technology can achieve, and it exposes the two limitations that the system proposed in this paper is designed to address: a heavy dependence on fossil energy, and a model accessible only to well-capitalised commercial operators. This section treats the Dutch sector as evidence on both counts, drawing on peer-reviewed studies and official statistics rather than on the promotional framing that often accompanies it.

\subsection{What the technology demonstrably achieves}

By gross export value the Netherlands is the second-largest agricultural exporter in the world after the United States, with agricultural exports of roughly \texteuro129 billion in 2024 \cite{cbs2025exports}, achieved on approximately 1\% of the European Union's agricultural land \cite[p.~2]{vangrinsven2019dutch}. This headline figure requires one immediate qualification: a large share of it is re-export of imported commodities passing through Rotterdam and Schiphol, so that the Netherlands' \emph{net} agricultural export value is closer to one third of the gross figure \cite[p.~2]{vangrinsven2019dutch}. The underlying productivity is nonetheless substantial: the country produces, for example, more than 9\% of EU dairy output and houses 6--8\% of its pig, poultry and meat production from that same 1\% land share \cite[p.~2]{vangrinsven2019dutch}.

The greenhouse method also raises yield per unit of land, which is the mechanism by which it reduces the land footprint of a given quantity of produce. Measured open-field tomato yields are on the order of 6 to 10~kg/m\textsuperscript{2}, against substantially higher yields under glass \cite[p.~16]{ntinas2017tomato}, and the high-tech Dutch sector operates at the upper end of that range and well beyond it. This land-sparing effect is the most robust of CEA's resource advantages, because it follows directly from measured productivity rather than from any particular energy or water accounting.

The reduction in pesticide use is equally well documented and, unlike some efficiency claims, specific to greenhouses. Integrated pest management based on biological control (the deliberate use of predatory mites, parasitoids and other natural enemies in place of chemical sprays) is used on over 90\% of greenhouse tomato, cucumber and sweet pepper production in the Netherlands \cite[p.~1992]{chandler2011biopesticides}. The enclosed environment is what makes this possible: pests can be excluded and beneficial organisms retained in ways that open fields do not permit.

Water use efficiency is consistently better under glass than in the open field. In the measured European comparison cited above, greenhouse cultivation converted irrigation water into marketable fruit more efficiently than open-field cultivation across every scenario, and the literature it reviews reports requirements as low as 14~L/kg for a heated Dutch rockwool system, against roughly 29 to 36~L/kg for unheated greenhouses elsewhere \cite[p.~16]{ntinas2017tomato}. The advantage is real, but its magnitude depends heavily on the baseline and the system boundary, and the very large reductions sometimes quoted reflect best-case comparisons rather than typical practice.

\subsection{The energy and carbon penalty}

The efficiencies above come at an energy cost that the promotional literature tends to omit, and which is decisive for this paper's argument. Controlled environments are productive partly because they substitute purchased energy for the free sunlight, ambient warmth and rainfall that open fields use, and in a temperate climate that substitution is expensive. Measured life-cycle studies bear this out. Across greenhouse and open-field tomato systems in Southern and Central Europe, the carbon footprint of greenhouse production ranged from about 0.4 to 7.2~kg CO\textsubscript{2}eq per kilogram of tomato, with fossil fuel for heating and electricity identified as the dominant contributor and open-field cultivation showing lower emissions and lower energy demand throughout \cite[Abstract, pp.~13--14]{ntinas2017tomato}. In heated Mediterranean greenhouses specifically, energy for heating alone accounted for 56.8\% of on-farm energy use, and heating raised emissions per kilogram by a factor of roughly 2.75 relative to unheated production \cite[Abstract]{perezneira2018tomato}. The penalty is therefore real and large, and it is driven by the source of the energy rather than by the method as such.

That distinction is what makes the penalty a design problem rather than a verdict. The same measured comparison found that when greenhouse production was supplied with renewable energy, well insulated, and run at higher marketable yield, its greenhouse-gas emissions fell sharply, in the best case to roughly 0.4~kg CO\textsubscript{2}eq per kilogram, comparable to open-field cultivation \cite[Abstract, p.~14]{ntinas2017tomato}. The carbon disadvantage of controlled-environment production is thus contingent on fossil energy, not intrinsic to the enclosure. This is precisely why the source of the energy, and the efficiency with which it is used, are central design variables for the system proposed in this paper.

The composition of that penalty points directly at the central design variable. A Dutch life-cycle assessment of lettuce found that vertical farming, which relies entirely on artificial lighting, carried a carbon footprint 5.6 to 16.7 times that of conventional production, with electricity for lighting and climate control accounting for about 85\% of the total \cite[Abstract]{blom2022embodied}. Lighting is thus the single largest determinant of how much energy a controlled environment consumes. A system that can draw on natural light, where the grower has the space and orientation to use it, faces a much smaller energy burden than one that must supply light artificially. The artificially lit mode is not thereby ruled out, since it extends production into seasons and locations that daylight alone cannot serve, but it does so at an energy cost that has to be measured rather than assumed. How far that cost can be met from on-site renewable generation, in either mode, is an open question that the system proposed in this paper treats as a primary design variable rather than a settled choice.

The Dutch sector itself illustrates how unresolved this problem remains. In 2023 it consumed about 92~PJ of energy and emitted approximately 5.1~Mt of CO\textsubscript{2}, with renewable sources supplying only 14\% of its energy, a share that fell for the first time on record that year and sits below the national average \cite[Summary; Sec.~2.5.2]{smit2024energiemonitor}. Most of its energy is natural gas burned in combined heat and power units, a technology that diffused rapidly through the sector after the Dutch energy market was liberalised in the early 2000s \cite[pp.~8--9]{vanderveen2015chp}; the sector now generates around 9.6~TWh of electricity, roughly 8\% of national consumption, and remains a net supplier to the grid \cite[Sec.~4.3]{smit2024energiemonitor}. On-site solar generation is poorly matched to greenhouses here: the roof has to transmit light to the crop rather than carry panels, and the electricity demand of supplementary lighting peaks when sunlight is scarce, in winter and at night \cite[Sec.~3.2.2]{smit2024energiemonitor}. Far from being incidental, fossil energy is built into the model as currently operated, and the sector acknowledges as much: under its 2022 energy-transition covenant it has committed to cutting residual emissions to 4.3--4.8~Mt CO\textsubscript{2}-equivalent by 2030 and to becoming climate-neutral only by 2040 \cite[Art.~1]{convenant2022glastuinbouw}. Decarbonisation is a stated ambition, not an achieved fact.

This is consistent with the broader environmental record of Dutch agriculture, which is not the unqualified success its export figures suggest. Including horticulture, the sector contributes about 14\% (32~Mt CO\textsubscript{2}-eq in 2016) of national greenhouse-gas emissions, and its emissions of ammonia and its surpluses of nitrogen and phosphorus per hectare are among the highest in the European Union \cite[Abstract]{vangrinsven2019dutch}. High per-product efficiency has coincided with high per-hectare environmental load, because the model concentrates production and inputs at extreme density.

\subsection{Structural limits: capital and the intensive model}

Two further limitations bear on whether the Dutch achievement can be generalised. The first is that it is expensive and institutionally specific. The sector's performance rests on high capital intensity, advanced proprietary technology, and an unusually well-resourced ecosystem of research, finance and specialised supply \cite[p.~2]{vangrinsven2019dutch}. These are precisely the conditions that smallholders, low-income communities, and regions without comparable institutions cannot reproduce. The efficiency is real; the access to it is not.

The second is that the intensive model has begun to reach hard ecological and legal limits. In 2019 a Dutch high court suspended permits for activities emitting nitrogen near protected nature reserves, stalling an estimated \texteuro14 billion in projects and prompting the government to order the world's densest livestock sector to contract \cite[p.~1180]{stokstad2019nitrogen}. This crisis is driven by livestock and manure rather than by greenhouses, so it does not bear on CEA directly; its relevance is as evidence that maximising output per hectare without regard to absolute environmental load is a strategy with a ceiling, one the Netherlands has now hit.

\subsection{Relevance to the proposed system}

Read in full, the Dutch case supports a precise and limited claim. It is an existence proof that controlled-environment cultivation delivers large, sustained reductions in land use and pesticide use, together with a real if more modest reduction in water use, for high-value produce. The technical core of the approach (closed-loop nutrient management, sensor-guided climate control, biological pest control) is demonstrably effective and is not in question.

What the Dutch case does not show is that this core, as currently packaged, solves the problems documented in Section~\ref{sec:problems}. On the contrary, it embodies two of them. Its carbon footprint is high and fossil-dependent, which means a CEA system can only contribute to the climate goals discussed in Section~\ref{sec:ghg} if it is paired with low-carbon energy and sited and crop-matched to keep the energy gap small, treating the choice between natural and supplementary lighting as a primary determinant of that gap rather than a settled preference. And its benefits are locked inside a capital- and institution-intensive commercial model, which is the inaccessibility that an open, decentralised design is meant to overcome. The system proposed in this paper takes the proven technical core of the Dutch approach and asks how to decarbonise it and open it: how to deliver the demonstrated resource efficiencies through a design that communities without concentrated capital can build, operate, and improve in common.
